\section{The Sovereign Second Brain and AION Flow}
\label{sec:sovereign-second-brain-flow}

\subsection{Product position}

Pilot is the user-facing surface, while AION is the durable, customer-controlled
second brain.  A foundation model is neither the product nor the owner of the
customer's accumulated intelligence.  Models, retrieval engines, agent harnesses,
cloud machines and business connectors are replaceable cognitive or operational
components used by AION for a bounded assignment.

The central product promise is therefore:

\begin{quote}
The customer keeps the brain and chooses the compute.  A provider can be removed,
a model can be replaced and a deployment can move, while identity, memory, the
business map, policy, capability authority and verified history remain intact.
\end{quote}

This separation permits a household, sole trader or small company to begin with a
local model and later attach private GPU compute.  A larger organisation may place
both AION and its models inside its own virtual private cloud.  A customer may also
connect an optional external API, provided that the disclosure is visible and
permitted before any information leaves the sovereign boundary.

\subsection{The three-layer architecture}

The product is divided into three stable layers:

\begin{enumerate}
  \item \textbf{AION brain.} Owns customer identity, scoped memory, the business
  map, policies, permissions, capability definitions, route learning and receipts.
  \item \textbf{Replaceable specialists.} Local or remote models, evidence
  systems, harnesses and compute targets perform language, vision, planning,
  retrieval or domain-specific calculations under a declared contract.
  \item \textbf{Pilot surfaces.} Phone, television, desktop, Workspace and
  Boardroom provide conversation, presentation, private review, approval and
  control.  They are surfaces of the same brain rather than independent brains.
\end{enumerate}

The canonical stores are provider independent: identity, memory, business map,
policy, capability and receipt.  A model adapter is forbidden from granting
authority, writing canonical identity or memory, promoting an inference directly
into the canonical business map, changing policy or erasing audit evidence.

\subsection{AION Flow: visual composition of intelligence}

The existing Glyph-based workflow canvas can be extended into \emph{AION Flow}, a
plug-and-play visual fabric for composing models, harnesses, evidence, compute and
authorized actions.  The useful form of the idea is not a simple chain in which
AION is only the first or last node.  AION must enclose the complete graph:

\begin{center}
\textbf{AION admission} $\rightarrow$ evidence/context $\rightarrow$
specialists $\rightarrow$ bounded comparison $\rightarrow$
\textbf{AION policy and approval} $\rightarrow$ capability $\rightarrow$
verification $\rightarrow$ \textbf{AION receipt and learning}.
\end{center}

At admission, AION resolves the active person, organisation, purpose, data scope
and minimum necessary context.  At the return boundary, it validates the proposed
result, checks evidence and authority, obtains any required private approval,
executes through a named capability, verifies the external outcome and records a
compact receipt.  A model result is always a proposal; it is never execution
authority.

\subsection{Node catalogue}

The first provider-neutral catalogue contains the following node families:

\begin{itemize}
  \item AION ingress, identity, context and memory-query nodes;
  \item local or authorized evidence and retrieval nodes;
  \item model nodes declaring exact model, provider, version and location;
  \item harness nodes defining prompts, tools, schemas and domain constraints;
  \item compute nodes selecting local hardware, private cloud or external APIs;
  \item comparison, critic, voting and refinement nodes;
  \item deterministic validator, contradiction and evidence-check nodes;
  \item policy, private approval and human-review nodes;
  \item governed capability and connector nodes;
  \item verification, outcome-measurement, receipt and bounded-learning nodes.
\end{itemize}

Each external node declares its provider, data classification, disclosure
destination, residency and encryption requirement.  Each route carries quality,
latency, cost, energy and context budgets.  A hidden fallback is not permitted to
cross a privacy, residency or cost boundary.

\subsection{Multi-model consensus and refinement}

Routing through several models can materially improve an outcome when the models
perform different roles or offer genuinely independent expertise.  Merely asking
the same model repeatedly does not establish correctness, and agreement between
models is not evidence.  A strong route may use one model to produce candidates,
another to identify failure modes, a retrieval node to supply current evidence and
a deterministic validator to test factual or structural requirements.

Every consensus or refinement loop must declare:

\begin{itemize}
  \item a maximum number of iterations;
  \item time, cost and token or compute budgets;
  \item measurable exit and failure conditions;
  \item the evidence required to resolve disagreement;
  \item whether a human reviewer is the final arbiter; and
  \item what happens when reliable agreement cannot be reached.
\end{itemize}

These controls prevent runaway agent conversations, correlated model errors and
the false appearance of certainty.

\subsection{Compilation and execution boundary}

Drawing a graph does not run it.  A visual flow first compiles into a hashed,
non-executing manifest.  The compiler requires exactly one AION ingress and one
AION receipt boundary, rejects nodes outside that enclosure, rejects provider
authority over canonical state, requires encryption metadata on external edges and
requires policy before every capability.  A consequential capability additionally
requires an approval ancestor.

Only after compilation may the existing Workflow Capsule machinery perform dry
run, permission evaluation, review and eventual execution.  Execution receipts
store route facts and content hashes without retaining unnecessary raw prompts,
private results or provider credentials.

\subsection{Portable brain contracts}

Version-one contracts now cover Brain Identity, Brain Export, Business Map and
Intelligence Route Receipt.  A complete Brain Export describes all six canonical
stores by schema version and content hash while explicitly excluding provider
credentials and model weights.  Business-map facts require provenance, confidence
and review state.  This allows a brain to move between a local machine, customer
server, private cloud or enterprise environment without changing its logical
identity.

\subsection{Public explanation surface}

Pilot Home now links to a public ``How the second brain works'' page.  It explains
the ownership model, deployment choices, governed route and AION Flow concept in
non-technical language.  The page is served locally without external scripts,
trackers or provider dependencies and may later be moved to the public Tessaris
website without changing the underlying product claim.

\subsection{Implementation and verification}

Phase zero of the Sovereign Business Brain programme is closed by the following
implemented artefacts:

\begin{itemize}
  \item provider-independent brain contracts and portable schema bundle;
  \item a documented reuse, adaptation and retirement inventory for existing
  COMDEX components;
  \item a versioned sovereign-boundary architecture decision and migration policy;
  \item a provider-free brain bootstrap with six canonical stores;
  \item a non-executing Sovereign Flow compiler and reference multi-model graph;
  \item tests for provider removal, restart, model replacement, forbidden provider
  authority, bounded consensus and governed capability execution; and
  \item the public Pilot dashboard explanation route.
\end{itemize}

The next build stage is the downloadable empty brain: signed installation,
hardware profiling, owner and mother-node bootstrap, encrypted vault, clean stores,
deployment choice, recovery, update, backup, restore and offline operation.
